\subsection{Modelo de programación de Sistemas Operativos}

\begin{frame}{Modelo de Programación de Sistemas}
 \begin{textblock*}{\textwidth}(0.4cm,1.2cm)
  \begin{itemize}[<+(1)->] \justifying \parskip 0.3cm
   \item Hasta ahora hemos trabajado con la vista de una arquitectura basada en el desarrollo de aplicaciones
   \begin{itemize} \justifying \parskip 0.3cm
    \item [\checkmark] Programación en C, C++, o Assembler en entornos Linux, o Windows.
    \item [\checkmark] Desarrollo de aplicaciones orientadas a objetos en vistosos entornos gráficos.
    \item [\checkmark] Desarrollo de aplicaciones en microcontroladores.
   \end{itemize}
   \item Trabajamos con un conjunto de recursos que se conocen como \textquotedblleft Modelo de programación de Aplicaciones\textquotedblright.
   \item Este modelo describe los registros de propósito general, las instrucciones, los tipo de datos, los modos de direccionamiento\dots tal parece que estaría todo dicho.
   \item Esto no es todo
  \end{itemize}
 \end{textblock*}
\end{frame}

\begin{frame}{¿Acaso hay mas recursos de desarrollo?}
 \begin{textblock*}{\textwidth}(0.4cm,1.2cm)
  \begin{block}<1->{Un paso mas hacia el hardware}
 Si nuestro requerimiento es diseñar un sistema que permita almacenar simultáneamente en memoria varias tareas, administrar su ejecución, evitar que se interfieran, y regular su acceso a los recursos de hardware disponibles, evidentemente estamos entrando en un mundo en el que la visión de Programador de Aplicaciones claramente no es suficiente.
  \end{block}
 \end{textblock*}
\end{frame}

\begin{frame}{¿Acaso hay mas recursos de desarrollo?}
 \begin{textblock*}{\textwidth}(0.4cm,1.2cm)
  \begin{itemize}[<+(1)->] \justifying \parskip 0.3cm
   \item Entonces se trata de desarrollar un conjunto de tareas independientes entre si que contribuyan a implementar el sistema requerido, y un software que las administre.
   \item Necesitamos contar con un SoC cuya CPU tenga capacidades para realizar el propósito planteado.
   \item Si nos las tiene, se puede desarrollar un sistema como el que nos estamos planteando, de todos modos. Pero no esperemos milagros en cuanto a rendimiento, ni le exijamos demasiado al modelo implementado.
   \item Primer paso: tener claro que CPU necesitamos para el requerimiento de nuestro sistema
  \end{itemize}
\end{textblock*}
\end{frame}

\begin{frame}{Modelo de Programación de Sistemas}
 \begin{textblock*}{\textwidth}(0.4cm,1.2cm)
  \begin{itemize}[<+(1)->] \justifying \parskip 0.3cm
   \item El modelo de programación de Aplicaciones de una CPU tiene una visión restringida del sistema en su conjunto.
   \item El modelo de programación de Sistemas, incluye un conjunto de recursos de hardware que amplía la capacidad de la CPU de resolver funciones que le permitan a un Sistema Operativo proveer los servicios y el entorno de programación que \textquotedblleft ven \textquotedblright las aplicaciones, de manera mucho mas veloz y eficiente.
   \item De hecho, el manejo de las excepciones y de las interrupciones de hardware que hemos visto anteriormente es claramente el terreno del Sistema Operativo.
   \item Y hay muchas mas funciones para proveer a un Sistema Operativo desde el hardware que no son necesarias para las aplicaciones.
  \end{itemize}
 \end{textblock*}
\end{frame}

\begin{frame}{Scheduling de Tareas}
 \begin{textblock*}{\textwidth}(0.4cm,1.2cm)
  \begin{itemize}[<+(1)->] \justifying \parskip 0.3cm
   \item Cualquier Sistema Operativo posee un módulo llamado \textbf{scheduler}, que le permite administrar la ejecución de tareas / procesos.
   \item Podemos considerar al \textbf{scheduler} dividido en dos capas: 
   \begin{itemize} \justifying \parskip 0.3cm
    \item [\checkmark]Una de alto nivel en la que se implementan una o mas \textit{políticas} para manejar el cambio de tareas, y que se traducen en diferentes \textit{algoritmos de scheduling}.
    \item [\checkmark]Una de bajo nivel en la que se realiza el \textit{context switch}.
   \end{itemize}
   \item En este punto nos vamos a concentrar en el módulo de \textit{context switch}, dejando para mas adelante los algoritmos de scheduling.
   \item El \textit{context switch} como veremos es físicamente dependiente del procesador y se escribe en assembler.
  \end{itemize}
 \end{textblock*}
\end{frame}

\begin{frame}{Capa de bajo nivel: Context Switch}
 \begin{textblock*}{\textwidth}(0.4cm,1.2cm)
  \begin{itemize}[<+(1)->] \justifying \parskip 0.3cm
   \item Un Context Switch genérico debe llevar a cabo las siguientes acciones:
   \begin{enumerate} \justifying \parskip 0.3cm
    \item Resguardar los registros core y cualquier otro del modelo de programador de aplicaciones en un área de memoria accesible solo para el kernel.
    \item En general el contexto forma parte de una estructura mas amplia llamada genéricamente \textbf{TCB} (por \textbf{T}ask \textbf{C}ontrol \textbf{B}lock) o en la jerga de los sistemas operativos de propósito general PCB (\textbf{P}rocess en lugar de \textbf{T}ask).
    \item Determinar la ubicación del TCB de la siguiente tarea en la lista de tareas en ejecución.
    \item Extraer el contexto de la nueva tarea y copiar registro a registro en la CPU.
   \end{enumerate}
  \end{itemize}
 \end{textblock*}
\end{frame}

\begin{frame}{Capa de alto nivel: Políticas de Scheduling}
 \begin{textblock*}{\textwidth}(0.4cm,1.2cm)
  \begin{itemize}[<+(1)->] \justifying  \parskip 0.3cm
   \item La capa de alto nivel llama al context switch cuando necesita conmutar tareas, en base a criterios definidos en su algoritmo.
   \item Antes de invocar al context switch evalúa una serie de variables:
   \begin{enumerate}
    \item prioridad,
    \item Preemption, 
    \item Atomicidad de operaciones, 
    \item concurrencia, 
    \item limitaciones en tiempo (latency), 
    \item paralelismo (en caso de que el sistema cuente con mas de una CPU), entre otros.
   \end{enumerate}
   \item Por lo tanto un \textbf{scheduler} resulta un código que puede oscilar entre algo muy sencillo, casi trivial, hasta algoritmos muy complejos que puedan considerar los aspectos señalados en el ítem anterior.
  \end{itemize} 
 \end{textblock*}
\end{frame}

\begin{frame}{Conclusiones}
 \begin{textblock*}{\textwidth}(0.4cm,1.2cm)
  \begin{block}<1->{Approach bootom-up} \justifying
  Intentaremos abordar el diseño de un scheduler mas o menos sencillo comenzando por la parte basal: El Context Switch.\\
  Esta es la parte dependiente del hardware. La capa superior interfaz consistente mediante puede reutilizarse, para lo cual es conveniente escribirla en C, para segurar su portabilidad.
  \end{block}
 \end{textblock*}
\end{frame}
